\section{Runtime and real-time operating points}
The runtime experiment measures scene analysis, state update, detector execution
at the selected resolution, confidence filtering, and tracker update on a Tesla
T4 over 200 frames. This distinguishes a controller that shapes compute over
time from a uniformly larger detector budget.

For ByteTrack, AC--MOT v3 has mean latency 74.3238 ms, P95 latency 100.6847 ms, and
throughput 13.4546 FPS; detector, tracker, and controller components are
54.5688, 7.9712, and 4.2562 ms. For OATrack, mean latency is 65.4748 ms, P95
85.8969 ms, and throughput 15.2731 FPS; the components are 52.6399, 1.8355,
and 4.2010 ms. The detector remains dominant and the controller is small but
nonzero, at roughly four milliseconds per frame.

The high-resolution baseline harness reports 139.93 ms and 7.15 FPS for
ByteTrack and 128.85 ms and 7.76 FPS for OATrack. Image-read work is not
included identically, so this is not a perfectly component-matched speedup.
The safe interpretation is reduced measured detector time from the selected
actions plus explicit controller overhead. Neither 13.4546 nor 15.2731 FPS is
claimed as 30-FPS real-time operation; the P95 values also do not establish a
hard per-frame deadline.

The tracker component differs substantially between hosts while detector times
are similar, supporting separate tracker-dependent reporting. A future
cost-matched study would need equal mean detector work and identical image-read
instrumentation to isolate scheduling from average compute.
